\subsubsection{Intraday Power Floating Leg}
\label{ss:intradaypowerfloatingleg}

An intraday power floating leg is specified in a \lstinline!LegData! node with \lstinline!LegType! set to \lstinline!IntradayPowerFloating!.
It is used to define a sequence of cashflows linked to an intraday power index over each calculation period.
Pricing is done in two aggregation steps: for each pricing date, a daily price per MWh is computed using the given intraday load shape; daily prices are then averaged over pricing dates in the period using daily volume weights.
The period amount is obtained by applying quantity, gearing and spread. The quantity interpretation depends on \lstinline!QuantityMode!: with \emph{LoadShapeMultiplier}, the load shape defines the absolute intraday and daily loads and
 \lstinline!Quantity! scales that absolute shape; with \emph{TotalEnergy}, the load shape is normalized and used as relative weights so that \lstinline!Quantity! is the total period energy.

For a detailed explanation of the pricing methodology, including the three key components (commodity curve, shape factors, and load profile), formulas for daily and period-average prices, and daylight-saving adjustments, see Section \ref{ss:intraday_power_pricing}.

The outline of an intraday power floating leg is given in listing \ref{lst:intradaypowerfloatingleg}. It has the usual \lstinline!LegData! elements described in section \ref{ss:leg_data} and an \lstinline!IntradayPowerFloatingLegData! node that is described in section \ref{ss:intraday_power_floating_leg_data} below.

\begin{listing}[h!]
\begin{minted}[fontsize=\footnotesize]{xml}
<LegData>
  <LegType>IntradayPowerFloating</LegType>
  <Payer>...</Payer>
  <Currency>...</Currency>
  <PaymentConvention>...</PaymentConvention>
  <PaymentLag>...</PaymentLag>
  <PaymentCalendar>...</PaymentCalendar>
  <ScheduleData>
    ...
  </ScheduleData>
  <PaymentDates>
    <PaymentDate>...</PaymentDate>
  </PaymentDates>
  <IntradayPowerFloatingLegData>
    ...
  </IntradayPowerFloatingLegData>
</LegData>
\end{minted}
\caption{Intraday power floating leg outline.}
\label{lst:intradaypowerfloatingleg}
\end{listing}

\subsubsection{Intraday Power Floating Leg Data}
\label{ss:intraday_power_floating_leg_data}

The \lstinline!IntradayPowerFloatingLegData! node outline is shown in listing \ref{lst:intraday_power_floating_leg_data}. The meaning and allowable values for each node are as follows:

\begin{itemize}

\item
\lstinline!Name!: Identifier specifying the intraday power index referenced by the leg.

\item
\lstinline!Quantities!: This node specifies constant or period-varying quantities. Usage is analogous to \lstinline!Notionals! in section \ref{ss:leg_data}.

\item
\lstinline!Spreads! [Optional]: Optional additive spread applied to the reference price in each calculation period.

\item
\lstinline!Gearings! [Optional]: Optional multiplicative gearing applied to the reference price in each calculation period.

\item
\lstinline!PricingCalendar! [Optional]: Calendar used to generate pricing dates inside each calculation period.

\item
\lstinline!IncludePeriodStart! [Optional]: If \emph{true}, include the period start date in pricing dates.

Allowable values: \emph{true}, \emph{false}. Defaults to \emph{true} if omitted.

\item
\lstinline!IncludePeriodEnd! [Optional]: If \emph{true}, include the period end date in pricing dates.

Allowable values: \emph{true}, \emph{false}. Defaults to \emph{false} if omitted.

\item
\lstinline!BusinessDays! [Optional]: If \emph{true}, pricing dates are business days in \lstinline!PricingCalendar!. If \emph{false}, pricing dates are non-business days. 
Use NullCalendar with business days set to \emph{true} for all calendar days.

Allowable values: \emph{true}, \emph{false}. Defaults to \emph{true} if omitted.

\item
\lstinline!PowerLoadProfileReference! [Optional]: Reference to an intraday power load profile convention. If no profile is given, it assumes a constant load over 24h.

\item
\lstinline!PowerLoadProfileData! [Optional]: Inline load profile data. If both \lstinline!PowerLoadProfileReference! and \lstinline!PowerLoadProfileData! are provided, the reference takes precedence. If neither is given, it assumes a constant load over 24h everyday. The structure and semantics of \lstinline!PowerLoadProfileData! are described in Section \ref{sss:intraday_power_load_conventions}.

\item
\lstinline!FXIndex! [Optional]: FX index used to convert each daily fixing before averaging.

\item
\lstinline!AvgPricePrecision! [Optional]: Non-negative integer precision for rounding the averaged price.

\item
\lstinline!Tag! [Optional]: Optional identifier used to link this leg with other related legs.

\item
\lstinline!QuantityMode! [Optional]: Quantity interpretation mode. 

\begin{itemize}
\item \emph{TotalEnergy}: The \lstinline!Quantity! is interpreted as the total MWh for the full calculation period. There are relative load factors derived from the given load profile. The total MWh are delivered relatively to the load profile.
\item \emph{LoadShapeMultiplier}: The \lstinline!Quantity! is interpreted as a multiplier applied to the period load implied by \lstinline!PowerLoadProfileData! (or \lstinline!PowerLoadProfileReference!). In this mode, a load profile is required.
\end{itemize}

Allowable values: \emph{TotalEnergy}, \emph{LoadShapeMultiplier}. Defaults to \emph{TotalEnergy} if omitted.

\end{itemize}

\begin{listing}[h!]
\begin{minted}[fontsize=\footnotesize]{xml}
<IntradayPowerFloatingLegData>
  <Name>...</Name>
  <Quantities>
    <Quantity>...</Quantity>
  </Quantities>
  <Spreads>
    <Spread>...</Spread>
  </Spreads>
  <Gearings>
    <Gearing>...</Gearing>
  </Gearings>
  <PricingCalendar>...</PricingCalendar>
  <IncludePeriodStart>...</IncludePeriodStart>
  <IncludePeriodEnd>...</IncludePeriodEnd>
  <BusinessDays>...</BusinessDays>
  <PowerLoadProfileReference>...</PowerLoadProfileReference>
  <PowerLoadProfileData>
    ...
  </PowerLoadProfileData>
  <FXIndex>...</FXIndex>
  <AvgPricePrecision>...</AvgPricePrecision>
  <Tag>...</Tag>
  <QuantityMode>...</QuantityMode>
</IntradayPowerFloatingLegData>
\end{minted}
\caption{Intraday power floating leg data outline.}
\label{lst:intraday_power_floating_leg_data}
\end{listing}
